\DocumentMetadata{
  pdfversion=2.0,
  pdfstandard=ua-2,
  lang=en-US,
  tagging=on,
  tagging-setup={math/setup=mathml-SE,tabsorder=structure}
}
\documentclass[11pt,letterpaper]{article}
\usepackage[margin=0.68in,includefoot]{geometry}
\usepackage{newcomputermodern}
\usepackage{amsmath,amscd,mathtools}
\usepackage{tikz,adjustbox}
\providecommand{\square}{\mdlgwhtsquare}
\usepackage[hidelinks]{hyperref}
\hypersetup{
  pdftitle={Numerical Analysis PhD Exam, May 1998},
  pdfauthor={University of Florida Department of Mathematics},
  pdfsubject={Graduate mathematics examination},
  pdfdisplaydoctitle=true,
  bookmarksopen=true
}
\setcounter{secnumdepth}{0}
\setlength{\parindent}{0pt}
\setlength{\parskip}{0.28em}
\setlength{\emergencystretch}{3em}
\setlength{\leftmargini}{2.15em}
\setlength{\leftmarginii}{2.65em}
\setlength{\leftmarginiii}{3em}
\pagestyle{empty}
\raggedbottom
\allowdisplaybreaks
\NewTaggingSocketPlug{block/list/label}{examproblem}{%
  \tagstructbegin{tag=Lbl}\tagstructend%
  \tagstructbegin{tag=\LBody}%
  \tagstructbegin{tag=\ProblemHeadingTag,title-o={\ProblemHeadingText}}%
  \tagmcbegin{tag=\ProblemHeadingTag,actualtext={\ProblemHeadingText}}%
  #2%
  \tagmcend%
  \tagstructend%
}
\newenvironment{examproblems}[2]{%
  \def\ProblemHeadingTag{#2}%
  \AssignTaggingSocketPlug{block/list/label}{examproblem}%
  \begin{enumerate}%
  \setcounter{enumi}{#1}%
  \renewcommand{\labelenumi}{\textbf{\theenumi.}}%
  \setlength{\topsep}{0.35em}%
  \setlength{\partopsep}{0pt}%
  \setlength{\itemsep}{0.48em}%
  \setlength{\parsep}{0.18em}%
}{\end{enumerate}}
\newenvironment{examsubparts}{%
  \AssignTaggingSocketPlug{block/list/label}{default}%
  \begin{enumerate}%
  \setlength{\topsep}{0.18em}%
  \setlength{\partopsep}{0pt}%
  \setlength{\itemsep}{0.16em}%
  \setlength{\parsep}{0.1em}%
}{\end{enumerate}}
\newenvironment{examinstructions}{%
  \par\begingroup\itshape
}{%
  \par\endgroup\vspace{0.18em}
}
\makeatletter
\renewcommand\subsection{\@startsection{subsection}{2}{0pt}%
  {-1.25ex plus -.25ex minus -.1ex}%
  {0.55ex plus .1ex}%
  {\normalfont\large\bfseries}}
\renewcommand{\@maketitle}{%
  \newpage\null\vskip -1em%
  {\footnotesize\bfseries
    \href{https://www.ufl.edu/}{University of Florida}, \href{https://math.ufl.edu/}{Department of Mathematics}\par}%
  \vskip -0.5em%
  \tagstructbegin{tag=H1,title={Numerical Analysis PhD Exam, May 1998}}%
  \tagpdfparaOff%
  \tagmcbegin{tag=H1}%
    {\LARGE\bfseries\href{https://gma.math.ufl.edu/exams/numerical-analysis-phd/}{Numerical Analysis PhD Exam}, May 1998\par}%
  \tagmcend%
  \tagpdfparaOn%
  \tagstructend%
  \vskip 0.25em%
  \tagmcbegin{artifact}%
    \hrule height 1pt%
  \tagmcend%
  \vskip 1em%
}
\makeatother
\title{Numerical Analysis PhD Exam, May 1998}
\author{}
\date{}
\begin{document}
\maketitle
\thispagestyle{empty}
\begin{examinstructions}
Do any 8 of the following 10 problems.
\end{examinstructions}
\begin{examproblems}{0}{H2}
\def\ProblemHeadingText{Problem 1.}
\item
This problem has four parts.
\begin{examsubparts}
\item[\textnormal{(1)}]
Give a careful description of a Householder transformation.
\item[\textnormal{(2)}]
Describe how Householder transformations can be used to bidiagonalize a \(5\times 5\) matrix.
\item[\textnormal{(3)}]
Describe the \(QR\) Iteration Algorithm with shift~\(\mu\).
\item[\textnormal{(4)}]
Describe how the \(QR\) Iteration is used to compute the \(SVD\).
\end{examsubparts}
\def\ProblemHeadingText{Problem 2.}
\item
This problem has two parts.
\begin{examsubparts}
\item[\textnormal{(1)}]
Give a careful statement and proof of Gershgorin’s Circle Theorem.
\item[\textnormal{(2)}]
Determine whether or not the matrix
\[
A=\begin{pmatrix}
-3&1&1\\
-1&-4&2\\
-1&0&-2
\end{pmatrix}
\]
 has the property that each eigenvalue has negative real part.
\end{examsubparts}
\def\ProblemHeadingText{Problem 3.}
\item
Suppose it is known that every eigenvalue of a matrix~\(A\) is real. Suppose the power method with shifts is to be used to approximate the largest eigenvalue of~\(A\). Explain how Gershgorin’s method can be used to aid in the choice of reasonable shifts.
\def\ProblemHeadingText{Problem 4.}
\item
This problem has three parts.
\begin{examsubparts}
\item[\textnormal{(1)}]
Give a careful statement of the \(LU\) factorization theorem.
\item[\textnormal{(2)}]
Give a careful statement and proof of the Cholesky Factorization Theorem.
\item[\textnormal{(3)}]
Find (if possible) a Cholesky factorization of the matrix
\[
A=\begin{pmatrix}
2&4&6\\
4&7&8\\
6&8&5
\end{pmatrix}.
\]
\end{examsubparts}
\def\ProblemHeadingText{Problem 5.}
\item
This problem has two parts.
\begin{examsubparts}
\item[\textnormal{(1)}]
Give a careful statement of the Sturm interlacing property for symmetric matrices.
\item[\textnormal{(2)}]
Determine whether or not the matrix
\[
A=\begin{pmatrix}
1&1&0\\
1&1&1\\
0&1&2
\end{pmatrix}
\]
 has any negative eigenvalues.
\end{examsubparts}
\def\ProblemHeadingText{Problem 6.}
\item
This problem has four parts.
\begin{examsubparts}
\item[\textnormal{(1)}]
Give a careful statement of the Pythagorean Theorem property for clamped cubic splines.
\item[\textnormal{(2)}]
Give a careful statement of the integral minimization property for clamped cubic splines.
\item[\textnormal{(3)}]
If \(s_n(x)\) denotes the clamped cubic spline interpolant for a function \(f(x)\) at the knots \((x_0,y_0),(x_1,y_1),\ldots,(x_n,y_n)\), then give a precise statement for the error between \(s_n(x)\) and \(f(x)\).
\item[\textnormal{(4)}]
If \(s_n(x)\) denotes the clamped cubic spline interpolant for a function \(f(x)\) at the knots \((x_0,y_0),(x_1,y_1),\ldots,(x_n,y_n)\), then give a precise statement for the error between \(s_n''(x)\) and \(f''(x)\).
\end{examsubparts}
\def\ProblemHeadingText{Problem 7.}
\item
Let \(f(x)\) be a real valued function on~\(\mathbb{R}\) with a continuous 3rd derivative at each point. Assume also that \(f(x)\) has a simple root at \(x=p\).
\begin{examsubparts}
\item[\textnormal{(1)}]
Show: If Newton’s method is used to approximate~\(p\), then there is a \(\delta>0\) with the property that the convergence rate to~\(p\) is 2nd order on the interval \([p-\delta,p+\delta]\).
\item[\textnormal{(2)}]
Approximate a solution to the following nonlinear system using Newton’s method. (Please don’t iterate—just set up the procedure for the initial guess \(x_0=(1,1,1)^t\).)
\[
\begin{aligned}
3x_1-\cos(x_2x_3)-0.5&=0,\\
x_1^2-625x_2^2&=0,\\
e^{-x_1x_2}+20x_3+9&=0.
\end{aligned}
\]
\item[\textnormal{(3)}]
What happens to the iteration process if the initial guess is changed to \(x_0=(1,0,1)^t\)?
\end{examsubparts}
\def\ProblemHeadingText{Problem 8.}
\item
Consider the method \(y_{k+1}=y_{k-1}+\frac{h}{2}(3f_k+f_{k-1})\) for solving the initial value equation \(y'=f(x,y)\), \(y(0)=\alpha\).
\begin{examsubparts}
\item[\textnormal{(1)}]
If the method is applied to the equation \(y'=-10y\), \(y(0)=4\), then is there a step size~\(h\) such that the global discretization error converges to zero as \(k\to\infty\)? Explain your answer.
\item[\textnormal{(2)}]
Consider the equation \(y'=xy(y-2)\) where \(y(0)=2\). Is this an ill-posed or a well-posed problem? Explain your answer. (Note that the general solution is \(y(x)=\frac{2y_0}{y_0+(2-y_0)e^{x^2}}\), when \(y(0)=y_0\).)
\end{examsubparts}
\def\ProblemHeadingText{Problem 9.}
\item
This problem has three parts.
\begin{examsubparts}
\item[\textnormal{(1)}]
Give a careful statement of the Contraction Mapping Theorem.
\item[\textnormal{(2)}]
Outline a proof of the Contraction Mapping Theorem.
\item[\textnormal{(3)}]
Indicate how Aitken’s method can be used to accelerate convergence.
\end{examsubparts}
\def\ProblemHeadingText{Problem 10.}
\item
This problem has two parts.
\begin{examsubparts}
\item[\textnormal{(1)}]
Explain Romberg’s technique for numerical approximation of an integral.
\item[\textnormal{(2)}]
Give a brief justification for the method.
\end{examsubparts}
\end{examproblems}
\end{document}
